% ======================================================================
%  Section 4.4.4 --- Acknowledged limitation.  Corrected and extended.
%
%  The three paragraphs you have are sound and the asymmetry-of-errors
%  argument is the right defence. Two changes and two additions.
%
%  1. STRENGTHENING. "its rows in L_R are small" is asserted. After the
%     addition to 4.4.1 it is a consequence: the row norm factorises into
%     the local gradient magnitude and a geometric term, and G is a
%     scale-aggregated gradient magnitude. The section's whole defence
%     rests on this step, so it should cite the place where it is proved
%     rather than restate it as an expectation.
%
%  2. THIS SUBSECTION IS WHERE 4.4.2 SENDS THE READER for the ablation of
%     the measure -- "Section X records what happens when it is used" --
%     and it does not record it. The paragraph is added below.
%
%  3. A SECOND LIMITATION IS MISSING, and it is one the chapter measures.
%     Section 4.3.3 identifies the early rejection of valid measurements as
%     the defect of a structure-blind estimator; 4.4 proposes the remedy;
%     and 4.7.3 then reports 13.6% of valid measurements strongly
%     down-weighted at iteration 1 with the remedy active. A limitation the
%     chapter's own results demonstrate belongs in the section that
%     acknowledges limitations, not left for a reader to notice.
%
%  4. TITLE. With three limitations the singular no longer fits, and 4.4.3
%     refers to this subsection as where "the scope of the method is
%     discussed". "Scope and limitations" matches both.
% ======================================================================

\subsection{Scope and limitations}
\label{sec:hermite-weight-limitation}

The structure measure~\eqref{eq:structG} reports that the image is locally
flat, not why it is flat. It therefore cannot distinguish a region that is
flat because it is occluded from one that is flat because the scene itself is
untextured there, and both are down-weighted alike. A genuinely untextured but
visible region will thus be suppressed even though its measurements are valid.

This behaviour is accepted, and is arguably appropriate rather than merely
tolerable. A region devoid of Hermite structure constrains the camera pose
only weakly, whatever the cause of its flatness. This is not an expectation
but a consequence of the factorisation established in
Section~\ref{sec:hermite-principle}: the norm of a row of $\mathbf{L}_{R}$ is
the product of the local gradient magnitude and a term depending only on the
image coordinates and the depth, and $G$ is a scale-aggregated gradient
magnitude, so a small $G$ entails small rows and a correspondingly slight
contribution to the descent direction. Down-weighting such a region therefore
discards little information in the case where the flatness is genuine, while
removing a substantial source of bias in the case where it is not. The
asymmetry of the two errors favours the rejection.
%%% (1) The cross-reference does the work that the sentence "its rows in
%%% L_R are small" previously did on its own.

%%% (2) The ablation 4.4.2 promises. Note the standard applied: a single
%%% run per variant is not grounds for preferring one on its final error,
%%% which is exactly the position 4.7.1 takes about the three variants.
%%% Applying it here too is what keeps the chapter consistent.
\paragraph{Ablation of the measure.}
An earlier formulation of~\eqref{eq:structG} used the summed Hermite
responses directly. Because those kernels carry a DC component --- their gain
at zero frequency equals their $L_{1}$ norm to three decimal places --- the
resulting measure reported local intensity rather than local structure, as
Figure~\ref{fig:structure-maps} shows and as the statistics quoted in
Section~\ref{sec:hermite-gate} quantify. That formulation attained the
smaller final error on the nominal run, $0.088^{\circ}$ against
$0.385^{\circ}$ in rotation. It is nevertheless not the measure adopted here,
for two reasons that do not depend on that comparison. The first is that the
argument of the preceding paragraph fails for it: a local average of the
intensity bears no relation to the norm of a row of $\mathbf{L}_{R}$, so the
weighting would down-weight dark regions rather than uninformative ones, and
would behave differently on a scene in which brightness and texture were
uncorrelated. The second is that it does not recover the identity when there
is nothing to reject: at convergence on a clean target it still rejects
$1.3\%$ of measurements and holds $8.7\%$ below a quarter weight, against
$0.3\%$ and $3.2\%$ for the measure adopted, because brightness does not
change as the pose converges. As to the difference in final error, a single
run of each variant is not a basis on which to prefer one, for the reason
given in Section~\ref{sec:nominal}; the choice rests on the two properties
above.

%%% (3) The limitation the results demonstrate.
\paragraph{The early rejection of valid measurements is not removed.}
Section~\ref{sec:tukey-limitation} identifies the systematic rejection of
valid measurements early in the servo as the characteristic defect of a
structure-blind estimator, and the weighting developed here is intended to
address it. It does not eliminate it. Section~\ref{sec:weight-maps} reports
that on an unoccluded target, with the structure normalisation active,
$13.6\%$ of measurements still carry a weight below a quarter at the first
iteration and $3.5\%$ are rejected outright, every one of them valid. The
normalisation rescales the residuals relative to one another, but while the
pose error is large the residual is large almost everywhere and the
distribution remains far from the one the estimator presumes. Whether the
normalisation reduces this over-rejection relative to a purely residual-based
estimator is not established here, no weight map having been recorded for
Variant~2, and is identified as an open point in
Section~\ref{sec:future}.
%%% This is the honest form. We can say the gate does not remove the
%%% problem, because Table 4.3 measures Variant 3. We cannot say it fails
%%% to improve on Variant 2, because Variant 2's map was never produced.
%%% If you run it, this paragraph gets a number and becomes a result.

It should nonetheless be recorded that the first of these limitations is a
real restriction of the method, and that a scene consisting predominantly of
untextured surfaces would defeat it, as it would defeat any scheme that
identifies unreliable measurements by their lack of local structure.

% ======================================================================
%  CHECK FOR REPETITION WITH 4.4.3
%  The corrected 4.4.3 now ends by saying that uniformity and not brightness
%  produces the vanishing response, and that a textured occluder would not
%  be suppressed. The first paragraph here says the measure cannot tell why
%  a region is flat. These are complementary, not duplicated -- 4.4.3 is
%  about bright versus dark, this is about occluded versus untextured -- but
%  read the two together once to be sure the transition does not repeat the
%  word "uniform" three times in four sentences.
%
%  LABELS USED ABOVE
%    sec:hermite-principle   4.4.1
%    sec:hermite-gate        4.4.2
%    sec:tukey-limitation    4.3.3
%    sec:nominal             4.7.1
%    sec:weight-maps         4.7.3
%    sec:future              future work
%    fig:structure-maps      Figure 4.1
%
%  IF YOU WOULD RATHER NOT CARRY THE SECOND ADDITION HERE, the alternative
%  is to state it in 4.7.3, where the measurement is. I recommend here: a
%  limitation stated in the methods section and then measured reads as
%  candour, whereas the same limitation appearing only in the results reads
%  as something the author discovered late.
% ======================================================================
